\chapter{Formy różniczkowe i całkowanie form}
\label{ma2:ch:differential-forms}

W poprzednich rozdziałach poznaliśmy dwie odmienne
operacje: całkowanie skalarnej gęstości względem
miary Hausdorffa oraz całkowanie pola wzdłuż
zorientowanej drogi. Pierwsza nie zależy od orientacji,
druga zmienia znak po odwróceniu kierunku.
Formy różniczkowe uogólniają drugą operację
na podrozmaitości dowolnego wymiaru.
Nie potrzebujemy ogólnej teorii rozmaitości
riemannowskich; pracujemy na podrozmaitościach
przestrzeni euklidesowej.

\section{Formy alternujące}

\begin{definicja}
\label{ma2:def:alternating-form}
Niech $V$ będzie rzeczywistą przestrzenią liniową
skończonego wymiaru. Formą alternującą stopnia $r\ge1$
nazywamy odwzorowanie $V^r\to\R$ liniowe względem
każdego argumentu i zerujące się, gdy dwa argumenty
są równe. Przestrzeń takich form oznaczamy
$\Lambda^r V^*$. Przyjmujemy $\Lambda^0 V^*=\R$.
\end{definicja}
Z wieloliniowości, po podstawieniu sumy dwóch
wektorów w dwie pozycje, wynika zmiana znaku
przy zamianie tych argumentów. Odwrotnie,
zmiana znaku daje zerowanie przy równych argumentach,
ponieważ pracujemy nad $\R$.

Dla funkcjonałów liniowych $\alpha_1,\ldots,\alpha_r$
określamy
\[
 (\alpha_1\wedge\cdots\wedge\alpha_r)(v_1,\ldots,v_r)
 =\det\bigl(\alpha_i(v_j)\bigr)_{i,j=1}^r.
\]
Wieloliniowość i alternowanie wynikają
z odpowiednich własności wyznacznika.
Dla $V=\R^2$ i funkcjonałów współrzędnościowych
$\mathrm dx,\mathrm dy$ wartość
$(\mathrm dx\wedge\mathrm dy)(v,w)$ jest
zorientowanym polem równoległoboku rozpiętego
przez $v,w$, a nie jego nieujemnym polem.
Zamiana kolejności wektorów zmienia znak.

\begin{figure}[htbp]
\centering
\begin{tikzpicture}[scale=1.05,>=Latex]
\begin{scope}
\fill[blue!15] (0,0)--(2,.35)--(2.7,1.8)--(.7,1.45)--cycle;
\draw[blue!65!black] (2,.35)--(2.7,1.8)--(.7,1.45);
\draw[->,thick] (0,0)--(2,.35) node[right]{$v$};
\draw[->,thick] (0,0)--(.7,1.45) node[above]{$w$};
\node at (1.35,-.65) {$(v,w)$: znak dodatni};
\end{scope}
\begin{scope}[xshift=5cm]
\fill[orange!20] (0,0)--(2,.35)--(2.7,1.8)--(.7,1.45)--cycle;
\draw[orange!80!black] (2,.35)--(2.7,1.8)--(.7,1.45);
\draw[->,thick] (0,0)--(2,.35) node[right]{$v$};
\draw[->,thick] (0,0)--(.7,1.45) node[above]{$w$};
\node at (1.35,-.65) {$(w,v)$: znak ujemny};
\end{scope}
\end{tikzpicture}
\caption{To samo pole geometryczne, przeciwne wartości formy.}
\label{ma2:fig:form-oriented-area}
\end{figure}

\begin{twierdzenie}[Baza form alternujących]
\label{ma2:thm:forms-basis}
Niech $(e_1,\ldots,e_m)$ będzie bazą $V$,
a $(e^1,\ldots,e^m)$ bazą do niej dualną, tj.
$e^i(e_j)=\delta_{ij}$.
Dla $1\le r\le m$ formy
$e^{i_1}\wedge\cdots\wedge e^{i_r}$,
gdzie $1\le i_1<\cdots<i_r\le m$, tworzą bazę
$\Lambda^r V^*$. W szczególności jej wymiar wynosi
$\binom mr$. Dla $r>m$ przestrzeń jest zerowa.
\end{twierdzenie}
\begin{proof}
Dla rosnących wielowskaźników $I,J$ wartości
formy bazowej o indeksie $I$ na wektorach
$(e_{j_1},\ldots,e_{j_r})$ są równe $\delta_{IJ}$,
co wynika z wyznacznika macierzy złożonej
z delt Kroneckera. Zatem formy są liniowo niezależne.
Dowolna forma alternująca $\alpha$ jest równa
\[
 \sum_{i_1<\cdots<i_r}\alpha(e_{i_1},\ldots,e_{i_r})
 e^{i_1}\wedge\cdots\wedge e^{i_r},
\]
gdyż obie strony mają te same wartości na
rosnących układach wektorów bazowych.
Na układach z powtórzeniem obie są zerowe,
a inne układy otrzymujemy przez permutację.
Wieloliniowość kończy dowód.
Dla $r>m$ każdy układ wektorów bazowych
ma powtórzenie, więc każda forma jest zerowa.
\end{proof}

\section{Iloczyn zewnętrzny}

Dla $\alpha\in\Lambda^r V^*$, $\beta\in\Lambda^s V^*$
definiujemy
\[
 (\alpha\wedge\beta)(v_1,\ldots,v_{r+s})
 =\frac{1}{r!s!}\sum_{\pi\in S_{r+s}}
 \operatorname{sgn}(\pi)
 \alpha(v_{\pi(1)},\ldots,v_{\pi(r)})
 \beta(v_{\pi(r+1)},\ldots,v_{\pi(r+s)}).
\]
Tutaj $S_\ell$ jest grupą permutacji $\ell$ indeksów.
Formy stopnia zero mnożymy jak skalary.
Wzór jest wieloliniowy, a permutacja argumentów
mnoży go przez znak tej permutacji (zmiana indeksu
sumowania), więc daje formę alternującą.

\begin{stwierdzenie}
\label{ma2:prop:wedge-algebra}
Iloczyn zewnętrzny jest dwuliniowy, łączny
i spełnia $\alpha\wedge\beta=(-1)^{rs}\beta\wedge\alpha$.
Dla form będących iloczynami funkcjonałów liniowych
jest zgodny z podanym wcześniej wzorem wyznacznikowym.
\end{stwierdzenie}
\begin{proof}
Dwuliniowość wynika wprost z definicji.
Dla iloczynów funkcjonałów liniowych rozwinięcie
obu wyznaczników w sumie daje pełne rozwinięcie
wyznacznika z $r+s$ wierszami: każda permutacja
występuje $r!s!$ razy z tym samym znakiem.
Stąd iloczyn dwóch takich form odpowiada
konkatenacji list funkcjonałów.
Łączność wynika dla nich z łączności konkatenacji,
a przestawienie dwóch bloków wierszy wymaga
$rs$ zamian, co daje czynnik $(-1)^{rs}$.
Twierdzenie \ref{ma2:thm:forms-basis} i dwuliniowość
rozszerzają obie własności na wszystkie formy.
\end{proof}

\begin{przyklad}
Dla $\alpha=\mathrm dx+2\mathrm dy$ i
$\beta=3\mathrm dx-\mathrm dy$ na $\R^2$
otrzymujemy $\alpha\wedge\beta=-7\mathrm dx\wedge
\mathrm dy$. Dla formy stopnia pierwszego
$\alpha\wedge\alpha=0$, ale dla form parzystego
stopnia nie musi tak być: na $\R^4$ forma
$\eta=\mathrm dx_1\wedge\mathrm dx_2+
\mathrm dx_3\wedge\mathrm dx_4$ spełnia
\[
 \eta\wedge\eta=2\mathrm dx_1\wedge\mathrm dx_2
 \wedge\mathrm dx_3\wedge\mathrm dx_4.
\]
Oba rachunki wynikają z dwuliniowości i reguły
zamiany czynników.
\end{przyklad}

\section{Formy różniczkowe i cofnięcie formy}

\begin{definicja}
\label{ma2:def:differential-form}
Na gładkiej podrozmaitości $M\subseteq\R^n$
formą różniczkową stopnia $r$ klasy $C^\ell$
nazywamy przyporządkowanie każdemu $p\in M$
elementu $\omega_p\in\Lambda^r(T_pM)^*$,
którego współczynniki w lokalnych parametryzacjach
są klasy $C^\ell$. Formy stopnia zero są funkcjami.
W przypadku podrozmaitości z brzegiem regularność
w parametryzacjach rozumiemy przez lokalne przedłużenia.
\end{definicja}

Na otwartym zbiorze $U\subseteq\R^m$ zapisujemy
jednoznacznie
\[
 \omega=\sum_{i_1<\cdots<i_r} a_{i_1\ldots i_r}(x)
 \,\mathrm dx_{i_1}\wedge\cdots\wedge\mathrm dx_{i_r}.
\]
Wynika to z twierdzenia \ref{ma2:thm:forms-basis}
w każdej przestrzeni stycznej.
Przykładem formy stopnia pierwszego jest różniczka
funkcji $f$, której wartością na wektorze $v$ jest
$Df(p)v$.

\begin{definicja}[Cofnięcie formy]
\label{ma2:def:form-pullback}
Niech $F:M\to N$ będzie odwzorowaniem klasy $C^1$.
Cofnięcie formy $\omega$ stopnia $r$ określamy przez
\[
 (F^*\omega)_p(v_1,\ldots,v_r)=
 \omega_{F(p)}(DF(p)v_1,\ldots,DF(p)v_r).
\]
Dla $r=0$ przyjmujemy $F^*f=f\circ F$.
\end{definicja}
Definicja różniczki na podrozmaitości pochodzi
z rozdziału \ref{ma2:ch:manifold-calculus}.
Dla form na zbiorach otwartych mamy
$F^*(\mathrm dy_i)=\mathrm dF_i$,
bo obie strony na $v$ dają $DF_i(p)v$.

\begin{stwierdzenie}
\label{ma2:prop:form-pullback}
Cofnięcie jest liniowe, zachowuje iloczyn zewnętrzny
oraz spełnia $(G\circ F)^*=F^*G^*$ i
$\id^*\omega=\omega$.
\end{stwierdzenie}
\begin{proof}
Liniowość wynika z podstawienia argumentów.
Zachowanie iloczynu zewnętrznego otrzymujemy
przez podstawienie $DF(p)v_i$ w sumie po
permutacjach definiującej ten iloczyn.
Reguła łańcuchowa dowodzi własności składania,
a różniczka identyczności jest identycznością.
\end{proof}
W szczególności zmiana lokalnej parametryzacji
przelicza współczynniki przez pochodne odwzorowania
przejścia. Jeżeli przejścia są gładkie,
klasa $C^\ell$ współczynników nie zależy od mapy.

\begin{przyklad}[Cofnięcie formy pola w biegunowych]
\label{ma2:ex:form-polar}
Dla $F(r,\theta)=(r\cos\theta,r\sin\theta)$
różniczkujemy obie składowe i otrzymujemy
\[
 \begin{aligned}
 F^*(\mathrm dx\wedge\mathrm dy)
 &= (\cos\theta\,\mathrm dr-r\sin\theta\,\mathrm d\theta)
 \wedge(\sin\theta\,\mathrm dr+r\cos\theta\,\mathrm d\theta)\\
 &=r\,\mathrm dr\wedge\mathrm d\theta.
 \end{aligned}
\]
Po zamianie kolejności parametrów forma ma
przeciwny znak, chociaż jakobian używany w całce
skalarnej nadal jest równy $r$ dla $r>0$.
\end{przyklad}

\section{Pochodna zewnętrzna}

\begin{definicja}
\label{ma2:def:exterior-derivative}
Dla formy $\omega=\sum_I a_I\,\mathrm dx_I$ klasy $C^1$
na otwartym zbiorze w $\R^m$, gdzie $I$ jest
rosnącym wielowskaźnikiem i $\mathrm dx_I$ oznacza
odpowiedni iloczyn zewnętrzny, określamy
\[
 \mathrm d\omega=\sum_{I,j}
 \frac{\partial a_I}{\partial x_j}
 \,\mathrm dx_j\wedge\mathrm dx_I .
\]
Dla form stopnia zero jest to zwykła różniczka.
\end{definicja}

\begin{twierdzenie}[Reguły pochodnej zewnętrznej]
\label{ma2:thm:exterior-derivative}
Dla form klasy $C^1$ pochodna zewnętrzna jest liniowa
i spełnia regułę
$\mathrm d(\alpha\wedge\beta)
=\mathrm d\alpha\wedge\beta+(-1)^r\alpha\wedge\mathrm d\beta$,
gdzie $r$ jest stopniem $\alpha$.
Jeżeli $\omega$ jest klasy $C^2$, to
$\mathrm d(\mathrm d\omega)=0$.
Dla $F$ klasy $C^2$ zachodzi
$\mathrm d(F^*\omega)=F^*(\mathrm d\omega)$.
\end{twierdzenie}
\begin{proof}
Liniowość wynika z definicji.
Dla $\alpha=a\,\mathrm dx_I$,
$\beta=b\,\mathrm dx_J$ stosujemy regułę iloczynu
dla funkcji $ab$. W drugim składniku przenosimy
$\mathrm db$ przez $r$ czynników $\mathrm dx_I$,
co daje znak $(-1)^r$. Sumowanie daje regułę ogólną.
Dla $\mathrm d^2\omega$ dostajemy sumę wyrazów
\[
 \frac{\partial^2 a_I}{\partial x_i\partial x_j}
 \,\mathrm dx_i\wedge\mathrm dx_j\wedge\mathrm dx_I.
\]
Wyrazy z $i=j$ znikają, pozostałe znoszą się parami
dzięki równości pochodnych mieszanych i alternowaniu.

W przypadku cofnięcia wystarczy rozważyć
$\omega=a\,\mathrm dy_{i_1}\wedge\cdots\wedge\mathrm dy_{i_r}$.
Jej cofnięcie to $(a\circ F)\,\mathrm dF_{i_1}\wedge
\cdots\wedge\mathrm dF_{i_r}$.
Z reguły iloczynu i $\mathrm d^2F_i=0$ wynika,
że pochodna zewnętrzna tego wyrażenia jest
$\mathrm d(a\circ F)\wedge\mathrm dF_{i_1}\wedge
\cdots\wedge\mathrm dF_{i_r}$.
Reguła łańcuchowa daje
$\mathrm d(a\circ F)=F^*(\mathrm da)$.
Własności z \ref{ma2:prop:form-pullback} kończą dowód.
\end{proof}

Na gładkiej podrozmaitości definiujemy $\mathrm d\omega$
lokalnie przez powyższy wzór w parametryzacjach.
Przemienność z cofnięciem zapewnia zgodność na
przecięciach map. W mapach brzegowych stosujemy
lokalne przedłużenia; ich pierwsze pochodne na brzegu
są wyznaczone przez ciągłość z wnętrza, więc wynik
nie zależy od przedłużenia.

\begin{przyklad}
Dla $\omega=P\,\mathrm dx+Q\,\mathrm dy$ na $\R^2$
mamy
\[
 \mathrm d\omega=(\partial_x Q-\partial_y P)
 \,\mathrm dx\wedge\mathrm dy.
\]
Wynika to z rozwinięcia dwóch różniczek
i $\mathrm dy\wedge\mathrm dx=-\mathrm dx\wedge\mathrm dy$.
Dla $\omega=-y\,\mathrm dx+x\,\mathrm dy$
otrzymujemy $2\,\mathrm dx\wedge\mathrm dy$.
W następnym rozdziale twierdzenie Stokesa połączy
całkę z tej formy po powierzchni z całką z $\omega$
po jej brzegu.
\end{przyklad}
